%%%%%%%%%%%%%%%%%%%%%%%%%%%%%%%%%%%%% 
% Check out the accompanying book, Even Better Books with LaTeX the Agile Way in 2023, for a discussion of the template and step-by-step instructions. https://amzn.to/3HqwgXM https://leanpub.com/eBBwLtAW/
% The template was originally created by Clemens Lode, LODE Publishing (www.lode.de), on 1/1/2023. Feel free to use this template for your book project! 
% I would be happy if you included a short mention in your book in order to help others to create their own books, too ("Book template based on \textit{Even Better Books with LaTeX the Agile Way in 2023} by Clemens Lode").
% Contact me at mail@lode.de if you need help with the template or are interested in our editing and publishing services.
% And don't forget to follow us on Instagram! https://www.instagram.com/lodepublishing/ https://www.instagram.com/betterbookswithlatex/
%%%%%%%%%%%%%%%%%%%%%%%%%%%%%%%%%%%%%

%%%%%%%%%%%%%%%
% Sections
%%%%%%%%%%%%%%%

% Subsections

% This is a comment about the next line.


% Use bold and capitals for the introductory sentence of a paragraph.
% use \textbf

% Use italics for emphasis only in rare cases, otherwise use bold.
% use \textit{}, \textbf

% Use quotes (like "To be or not to be")sparingly. When using quotes, use italics.

% use \textit{}

% References to titles: italics. Check the preamble in main.tex

% When referencing the name of your book, use bold, caps, and italics.
\newcommand{\bookname}[1]{\textbf{\textit{\uppercase{#1}}}}


% For titles, use bold, capitals, a different font, a different color, and initial caps.
\definecolor{dark-blue}{RGB}{28, 69, 135}
% \addtokomafont{section}{\color{dark-blue}\LARGE\sffamily\bfseries}
\addtokomafont{section}{\LARGE\sffamily\bfseries}
% \addtokomafont{subsection}{\color{purple}\large\sffamily\bfseries}
\addtokomafont{subsection}{\large\sffamily\bfseries}

% If you want to suppress chapter numbers from chapter and section titles for easier reading, uncomment the following lines.
%\renewcommand*\chapterformat{}
%\renewcommand{\thesection}{\arabic{section}}

\ifxetex
%-------------------------------------------
% Define space before and after a chapter, section, and subsection.
\renewcommand*\chapterheadstartvskip{\vspace*{-3\topskip}}
\renewcommand*\chapterheadendvskip{\vskip-.5\baselineskip\noindent{\color{gray}\rule{\linewidth}{2pt}}\par}
\newcommand\subsectionprelude{\vspace{-\parskip}}
\newcommand\subsectionpostlude{\vspace{-\parskip}}
\RedeclareSectionCommands[
    beforeskip=.5\baselineskip,
    afterskip=0.25\baselineskip
]{section,subsection,subsubsection}


% Use this code to define space before and after sections.
\makeatletter
\let\origsubsection\subsection
\renewcommand\subsection{\@ifstar{\starsubsection}{\nostarsubsection}}
\newcommand\nostarsubsection[2][\relax]{
  \subsectionprelude
  \ifx\relax#1\origsubsection{#2}\else\origsubsection[#1]{#2}\fi
  \subsectionpostlude}
\newcommand\starsubsection[2][\relax]{
  \subsectionprelude
  \ifx\relax#1\origsubsection*{#2}\else\origsubsection*[#1]{#2}\fi
  \subsectionpostlude}
\makeatother

\fi
% Set the space between paragraphs and deactivate indentation for paragraphs.
% \setlength{\parskip}{1.4\baselineskip}
\setlength{\parskip}{1.0\baselineskip}
\setlength{\parindent}{0pt}


%-------------------------------------------
% Use this command to reference literal names of real-world items that your book discusses (for example, flashcards or playing cards); it will appear in dark red and capitalized.
\newcommand{\cardref}[1]{{\textsc{\color{red}#1}}}

%-------------------------------------------
% Use this command to use small capitals in dark blue when referencing core concepts.
\newcommand{\concept}[1]{\textsc{\lowercase{\color{blue}#1}}}



%-------------------------------------------
% Reduce the space between each item of itemized lists.
\setlist{nosep}

% Reduce the space before and after a list (the itemize environment) to 5pt.
\setlist[itemize]{parsep=5pt}

%%%%%%%%%%%%%%%
% Lists
%%%%%%%%%%%%%%%

% \setlist[itemize,1]{wide = 0pt,labelwidth = 0cm, label=\color{red}\ding{110}} 
% \setlist[itemize,1]{wide = 0pt,labelwidth = 0cm, label=\ding{110}} 

% Use level 2 to denote a list within a list.
% \setlist[itemize,2]{wide=4pt,leftmargin=4pt, label=\color{red}\ding{121}} 
% \setlist[itemize,2]{wide=4pt,leftmargin=4pt, label=\ding{121}} 

%\setlist[itemize,3]{label={}} %list within a list within a list
%\setlist[itemize,4]{label={}} %list within a list within a list within a list

%%%%%%%%%%%%%%%
% Numbered lists
%%%%%%%%%%%%%%%

% Reduce the space between each item of enumerated lists.
\setenumerate{nosep}

% Format listings (grey background).
\lstset{breaklines=true,backgroundcolor=\color{lightgray},tabsize=1,basicstyle=\ttfamily\footnotesize}

%%%%%%%%%%%%%%%
% Page background
%%%%%%%%%%%%%%%

% Note that background pictures require the use of bleed for print (see /lib/bookformat.tex).
\ifxetex
\AddToHook{shipout/background}[jinwen/opac]{
    \put (0in,-\paperheight){\includegraphics[width=\paperwidth,height=\paperheight]{images/white.png}}
}




%%%%%%%%%%%%%%%%%
% Page background pictures
%%%%%%%%%%%%%%%%%

% Note that background pictures require the use of bleed for print (see /lib/bookformat.tex).

\newcommand{\cutlargepic}[1]{\AddToHookNext{shipout/background}{\put (0in,-22.8em){\includegraphics[width=\paperwidth]{#1}}}~\vspace{17em}}

\newcommand{\cutpic}[1]{\AddToHookNext{shipout/background}{\put (0in,-17.8em){\includegraphics[width=\paperwidth]{#1}}}~\vspace{14em}}
   
\newcommand{\cutbottompic}[1]{\AddToHookNext{shipout/background}{\put (0in,-\textheight){\includegraphics[width=\paperwidth]{#1}}}}    
\else

% You cannot use background pictures in e-books, so just add them as normal pictures or hide them.

\newcommand{\cutlargepic}[1]{\includegraphics[width=\paperwidth]{#1}}
\newcommand{\cutpic}[1]{\includegraphics[width=\paperwidth]{#1}}
\newcommand{\cutbottompic}[1]{\includegraphics[width=\paperwidth]{#1}}

\fi

% Use these commands to create custom symbols that can be used as part of the text.
\newcommand{\vcenteredinclude}[1]{\begingroup\setbox0=\hbox{\includegraphics[width=0.1in]{#1}}\parbox{\wd0}{\box0}\endgroup}
\newcommand{\squareIcon}{\vcenteredinclude{images/square.png}}


